\section{Method}

\smallskip\noindent\textbf{Backbones.} Two frozen text-to-music models sit behind one interface. Stable Audio Open 1.0 \citep{evans2024sao} is a latent diffusion transformer with a $v$-prediction objective, sampled with DPM-Solver++(2M) for 50 steps at guidance 7. ACE-Step 1.5 XL turbo \citep{gong2026acestep} is a rectified-flow transformer, $x_t=(1-t)x_0+t\epsilon$ with the network predicting $\epsilon-x_0$, distilled to 8 Euler steps without guidance. Everything below is written for a generic prediction $f_\theta(z_t,t,c)$.

\smallskip\noindent\textbf{A slider.} A slider is a set of low-rank updates, one per linear layer in the attention and feed-forward blocks. With position $s\in\mathbb{R}$ a layer computes
\begin{equation}
  y = Wx + s\,\frac{\alpha}{r}\,BAx ,
\end{equation}
with $r=4$ and $\alpha=1$. $s$ is read at every forward pass and can differ per sample. Sliders add, and a slider can be switched on only for $t\le t_0$ (\emph{gating}).

\smallskip\noindent\textbf{Training from a prompt pair.} Let $c^+$ and $c^-$ be a base prompt $c$ with a phrase for each end of the attribute appended. Following Concept Sliders, for $s\in\{+1,-1\}$ the slider is trained so that
\begin{equation}
\begin{split}
  f_{\theta,s}(z_t,t,c) \approx{}& f_\theta(z_t,t,c) \\ &+ s\,\eta\,\bigl(f_\theta(z_t,t,c^+)-f_\theta(z_t,t,c^-)\bigr),
\end{split}\label{eq:pair}
\end{equation}
with the right-hand side from the frozen model. $z_t$ comes from sampling for a random number of steps with the slider active. We use $\eta=4$ for Stable Audio Open and $\eta=2$ for ACE-Step, 1{,}000 iterations, and 48 training prompts disjoint from the evaluation prompts.

\smallskip\noindent\textbf{Training from two sets of clips.} The second trainer uses no text. Given sets $H$ and $L$ of latents, we minimise the model's denoising loss with the slider at $+1$ on $H$ and at $-1$ on $L$:
\begin{equation}
\begin{split}
  &\mathbb{E}_{x_0\in H}\,\bigl\|f_{\theta,+1}(z_t,t,c)-u\bigr\|^2 \\ &+ \mathbb{E}_{x_0\in L}\,\bigl\|f_{\theta,-1}(z_t,t,c)-u\bigr\|^2 ,
\end{split}
\end{equation}
where $u$ is the regression target for $(x_0,\epsilon,t)$. One update serves both ends with opposite sign, so to first order what the sets share cancels. $H$ and $L$ are the top and bottom 20--30\% of the model's own clips along some per-clip measurement, taken \emph{within each prompt} so that prompts are balanced. The measurement can be a descriptor, a quality score, or the projection on a direction in an embedding space, and it can be residualised against other measurements first. No differentiable path through the decoder or the embedding model is needed.

\smallskip\noindent\textbf{Axes.} We take directions from three places. (i) Principal components of CLAP embeddings of 1{,}024 clips of one concept, as in SliderSpace. (ii) Principal and independent components \citep{yamagiwa2023ica} of CLAP \citep{wu2023clap} and MuQ-MuLan \citep{zhu2025muq} embeddings of \covReal{} real recordings from FMA \citep{defferrard2017fma}, after removing the 40\% most vocal clips and each genre's mean. (iii) Principal components of the generator's own activations: the mean cross-attention output of every block for 768 clips, with each prompt's mean removed.

\smallskip\noindent\textbf{Other ways to move an attribute.} Under identical prompts, seeds, and sampler we also run the right-hand side of Eq.~\eqref{eq:pair} directly at every step (the training-free recipe of FreeSliders, three forward passes per step); activation steering, which adds the mean difference in cross-attention output between $c^+$ and $c^-$ to every block \citep{staniszewski2026tada}; interpolation of the text conditioning toward $c^+$ or $c^-$; and, where one exists, a conventional effect applied to the unsteered clip.
